% !PRINTED 743
\weberpage{743}
Man könnte dann an Stelle der Matrix \S.\,201, (20) eine Matrix
der $b_{k,i}$ von $\nu'+1$ Reihen setzen, und es würde sich auf
dem im vorigen Paragraphen eingeschlagenen Wege schliessen
lassen, dass $M$ durch $2^{\nu'+1}$ theilbar sein müsste, was
nicht der Fall ist.

Wenden wir aber das System der Multiplicatoren $b_i$, das den
Bedingungen (1) genügt, auf die Gleichungen (14) des vorigen
Paragraphen an, indem wir mit $b_i$ multipliciren und dann
summiren, so ergiebt sich, wenn wir Vielfache von 2 weglassen
und dies auch hier (wo irrationale Zahlen, nämlich die
Logarithmen, auftreten) durch das Zeichen der Congruenz
ausdrücken:
\weq{2}{
\begin{aligned}
l'_1 \equiv{} & L_0 \sum M_{0,i} b_i + L_1 \sum M_{1,i} b_i + \cdots \\
& + L_{\nu'-1} \sum M_{\nu'-1,i} b_i \pmod{2}.
\end{aligned}
}
Wollen wir die Congruenz in eine Gleichung verwandeln, so kommt
eine Summe von Grössen
\[ l_1, l_2, \ldots, l_{\nu-1}, \quad l'_1, l'_2, \ldots, l'_{\nu'-1} \]
mit \gesperrt{geraden ganzen Zahlen} als Co\"efficienten hinzu.

Dieselbe Betrachtung lässt sich aber auf $l'_2, l'_3, \ldots,
l'_{\nu'-1}$ anwenden, und danach kann man, wenn man von den
Logarithmen wieder zu den Zahlen übergeht, die Formel (2) so
in Worte fassen:

1. \gesperrt{Jede Einheit des Körpers $H_\mu$ ist als Product
einer Normaleinheit des Körpers $H_m$ und eines Quadrates
einer Einheit in $H_m$ darstellbar.}

Wenn wir darauf den Satz 2., \S.\,200 anwenden, so ergiebt
sich, dass jede Einheit $\mathcal{E}(r)$ des Körpers $H_\mu$,
die mit allen ihren Conjugirten dasselbe Zeichen hat, vom
Vorzeichen abgesehen, das Quadrat einer Einheit in $H_m$ sein
muss. Ist aber
\[ \pm\, \mathcal{E}(r) = \Delta(r)^2, \]
worin $\mathcal{E}(r)$ eine Einheit in $H_\mu$ und $\Delta(r)$
eine Einheit in $H_m$ ist, so muss auch $\Delta(r)$ in $H_\mu$
enthalten sein.

Denn setzen wir
\[ \Delta(r) = \Delta_1(r^2) + r \Delta_2(r^2), \]
so kann das Quadrat davon nur dann in $H_\mu$ enthalten sein,
also durch die Vorzeichenänderung von $r$ ungeändert bleiben,
wenn entweder $\Delta_1$ oder $\Delta_2 = 0$ ist. Es kann aber
nicht $\Delta(r) = r \Delta_2(r^2)$ sein, denn sonst wäre
$\Delta_2(r^2)$ eine Einheit in $\Omega_\mu$ und müsste
%CONT>>
